\chapter[Sermon 91. That the World Shall Be Renewed, the Saints Glorified and Made Blessed: And What That Felicity Shall Be, and How Certain.]{That the World Shall Be Renewed, the Saints Glorified and Made Blessed: And What That Felicity Shall Be, and How Certain.}
\label{sermon:091}
\markright{Sermon 91. That the World Shall Be Renewed, the Saints Glorified ...}

And I saw a new heaven and a new earth. For the first heaven and the first earth had vanished, and there was no more sea. And I, John, saw that holy city, New Jerusalem, coming down from God out of heaven, prepared as a bride adorned for her husband. And I heard a great voice from the throne, saying, “Behold, the tabernacle of God is with men, and he will dwell with them. They shall be his people, and God himself will be with them, and will be their God. And God will wipe away every tear from their eyes. There shall be no more death, neither sorrow, neither pain. For the former things have passed away.” And he who was seated on the throne said, “Behold, I am making all things new.” And he said to me, “Write this, for these words are true and faithful.” And he said to me, “It is done.”

\index{Antichrist}I warned you concerning the beginning of chapter 15 of this book, that a fifteenth part of this work begins at chapter 15 and deals with God’s righteous and just judgments. And since God’s judgments are of two kinds—when he punishes evil according to their wickedness, and rewards the good with rewards—I said that this passage consists of two parts. First, I said that John abundantly describes the torments to be inflicted on Antichrist and all the ungodly. Secondly, he describes rewards, especially at the end of the world, to be given to all the saints. For we have often heard in this book that the souls separated from the body are, immediately after bodily death, taken up into everlasting life, but that the greatest felicity of all awaits the faithful at the end of the world, when the bodies are raised again and receive the rewards of everlasting glory. This passage is treated throughout chapter 21 and the beginning of chapter 22. Just as in the earlier part he has shown hell as if wide open, and the everlasting torments as if they could be seen presently, so in this later part he unlocks, as it were, or opens heaven itself, so that with the eyes of faith, we may see what hope and glory awaits the saints. And with all, the article of our faith is most clearly explained: “I believe in everlasting life.” And again, for greater clarity, he declares these things through a vision, which others note as the seventh and last. Therefore, all things are figured spiritually, not carnally, to be understood and taken. Doubtless, the matters are excellent even when understood literally; however, we must think of spiritual matters, and always consider them greater than human speech can express. For we know, as taught by the doctrine of the Prophets and Apostles, that it is always true that what no eye has seen, nor ear has heard, nor has entered the heart of man, is what God has prepared for those who love him. 1 Corinthians 2

\index{Apocalypse (Book of)}\index{Ezekiel (Book of)}And the chief points of this passage are these. First, he shows that the world shall be renewed. Secondly, he signifies that the saints shall be glorified and blessed. And he declares in general what that blessedness shall be. Immediately thereafter, he confirms these things with many reasons; moreover, he describes the place, the court, and palace of the blessed, and likewise the glory and blessedness of the saints. Which in the beginning of chapter 20, he finishes exceedingly well, under the figure of a river and tree of life. And just as he has for the most part borrowed all his things out of the books of prophets, which John also with his Apocalypse illuminates, so has he at this present borrowed these from chapters 65 and 66 of Isaiah, and chapter 37 of Ezekiel, and the last chapters of the same.

\index{Augustine, Saint}\index{Arethas of Caesarea}\index{Primasius}Of the renewing of the world he speaks plainly, as does also the Apostle Peter in his second Epistle, chapter 3, that all things truly should be purged by fire, and not wholly abolished and annihilated, but should certainly be purified from all corruption: for Arethas says that he does not signify the extinguishing of the creation, but a renewing for the better. Therefore, John says expressly that he saw a new heaven and a new earth, and he adds by explanation that the first heaven and the first earth have vanished away: that is, they are changed in their qualities, so that the corruptible things are now gone, created for corruptible uses. For even so the sea is no more, also certainly subject to corruption, but changed into something better. Augustine and his colleague Primasius suppose that the troubled state of the world (signified not seldom in the scriptures by the sea) about the end of the world shall cease. Reading chapter 17 of the 20th book \textit{De civitate dei}, he reasons likewise at length about this innovation of the world, in the same 20th book \textit{De civitate dei}, and chapter 18 and other places. I think it best in this matter to put away all curiosity: and if any hidden thing appear therein, that it be reserved until that day, in which we shall see all things evidently. And I suppose that these things concerning the renewing of heaven and earth are not spoken so that there should any place be prepared for us, which we should inhabit again in these lower parts under heaven (for we believe that we shall fly up into heaven and go meet the Lord in the clouds, according to the doctrine of the Apostle, 1 Thessalonians 4), but so that our minds are thus confirmed, that the faithful shall undoubtedly be renewed and glorified. For if heaven and earth, made for man, be renewed and purified, who will doubt now that men themselves shall be most chiefly clarified?

Consequently, now John declares that the saints will both be renewed and glorified and placed in blessed seats, and he signifies generally what the glory of the saints will be. Afterward, he will declare more fully and individually all those things most diligently. For he hears an angel saying: “Come, I will show you the bride, the wife of the Lamb,” and so on. He now names this figuratively a city, and that city is truly holy, the new Jerusalem. And a city signifies both the place and habitation, as well as those who dwell in the place, that is, the citizens themselves. This city is therefore not only the place of the blessed, but also the very communion of saints, previously prefigured in the city of Jerusalem. But he puts a great difference between this new one and that visible, physical Jerusalem. For he calls ours holy; that other in the land of Palestine was profane, polluted with the blood of Christ, prophets, and apostles, and for the same reason utterly destroyed. Ours is also called new, for the communion of saints will be renewed at the same time. Therefore, it follows in interpretation, “coming down from heaven,” not that the habitation of the saints after judgment will again be on earth, but that the glory and renewal will be granted from heaven by divine majesty and power. As James is also reported to have said, every good gift and every perfect gift is from above, coming down from the Father of lights. And Paul also, in Galatians 4, said that the free church is the heavenly Jerusalem. The same in 1 Corinthians 15: The first man is of the earth, earthly; the second man is the Lord himself from heaven. Those who are earthly are like that earthly one; those who are heavenly are like that heavenly one. And just as we have borne the image of the earthly man, so we will also bear the image of the heavenly man. Therefore, John rightly said that the church of the saints comes down from heaven, that is, from heaven receiving its glory. For again, by demonstration, he says, “prepared by God, as a bride adorned for her husband.” The apostle, in 2 Corinthians 5, says, “We know that if our earthly house, this tabernacle, is destroyed, we have a building from God, a house not made with hands, eternal in heaven.” And immediately afterward: “He who has prepared us for this is God.” He removes all corruption from his saints, but gives and teaches them to be purified with all gifts of the body, so that they may be worthily adorned and dwell in the everlasting bridal chamber with their bridegroom, Christ. Therefore, this adorning consists in the abolishing of all corruption and mortality, and in the gift of incorruption, immortality, and glory. Paul, the apostle, also speaks of the purifying and adorning of the bride in Ephesians 5. In this world, the purging and trimming begin, and finally, at the end, it is finished most perfectly. For then the church will have neither spot nor wrinkle, all corruption truly wiped away, and all glory received. And here learn by the way that the saints are prepared by God; therefore, salvation is of mere grace.

And he proceeds to declare yet more plainly, what the glory shall be: a subject on which he has been occasioned to speak more often than once. Blessedness chiefly consists in two things. For God will give unto his Saints all that is good, and will take from them all evil: and so shall these for ever enjoy the ultimate good, and the most perfect felicity, and shall lack all pain and misery. Saint Augustine in the end of his book \textit{The City of God}: “How great,” he says, “shall that felicity be, where no evil shall be, no good shall lack?” And this declaration of eternal felicity has its parts, by which it is made manifest. For first, a voice, and that a great one, cried from the Throne: “Behold, the tabernacle of God is with men.” The conjunction of God with holy men was in time past prefigured by the Tabernacle of Witness, whereby God testified that he would be in the midst of his people. And the same he will at the end, after the judgment, fulfill most abundantly. And therefore that voice adds: “and he will dwell with them, and they shall be his people, and God himself with them, and will be their God.” The which Saint Paul seems to have uttered more succinctly and briefly, “and God shall be all in all.” For whatever is good, whatever is fair, whatever is pleasant and delectable, whatever the mind of man can imagine to be wished for, briefly, whatever pertains to the true and perfect felicity and blessed life, that same great God almighty will be wholly, and will show in himself most fully. And just as all and singular men enjoy to a pleasant satiety the amiable brightness and some heat of the sun, yet the sun nevertheless loses nothing by it: and although all men use the sun in common, each nevertheless enjoys it as proper and peculiar; right so in another world we shall use that eternal light, and joy everlasting and unspeakable. Whereof more plentiful things shall follow shortly.

And then, just as God in himself gives to the glorified all goodness, so will he remove all evil from them; so that they will not only be delivered from calamities, but those calamities shall never return, nor be feared again. This he declares most abundantly through words taken from the oracles of the prophets. “God will wipe away all tears from their eyes,” he says. He used this kind of speech also in chapter 7, truly taken from chapters 25 and 65 of Isaiah. And David also in Psalm 126: “They who sow in tears shall reap in gladness.” He seems to have alluded to mothers, who wipe the eyes of their tender children, comfort the sorrowful, and cherish them hurt or bruised. Therefore, if the saints have suffered any pain or grief in this world, it will be requited to them, provided that they shall feel no more adversity. The Lord also said in the gospel: “Verily, verily, I say unto you, you will weep and lament, but the world will rejoice; and you will be sorrowful, but your sorrow will be turned into joy,” and so on, in John 16.

\index{Resurrection}Consequently, he declares even more fully by naming the calamities, that the saints in the next life will be delivered at once from all evil: and death will be no more. For they will be rewarded with eternal life. Therefore, there will be no more fear of death, which is in a manner more bitter than death itself. The same affirms the Apostle in 1 Corinthians 15, citing the testimony of the prophet Hosea. There will be no mourning nor sorrow, which dries up the bones, even though they are full of vitality. For the joy of the saints will be perpetual. There will be no clamor, no complaints, no questioning or grumbling. For there will be no injustice, no malice, or envy. This world sounds and reverberates full of the clamor and cries of poor wretches. But in the blessed realms there will be no misery. There will be no pain, labor, sickness, or weariness. The cause of this is that the old things are gone. There is now another life, indeed, a very different manner of living from that which we live now. Therefore, whatever is of sin and subject to corruption will be taken away, as the Lord said in the Gospel, “The children of this world marry and are married: but those who are accounted worthy to attain to that other world and resurrection from the dead will neither marry nor be married. For they can die no more; for they are made like angels, and are the children of God, because they are the children of the resurrection.” Luke 20. But of eternal life we have spoken more in our commentaries on Matthew 12. And the Lord himself in John gathers the sum of all things and says how he makes all things new. Therefore, in the world to come we shall think of no carnal or corrupt thing, but a heavenly one.

But the minds of the faithful are grievously tempted in this matter, the Devil suggesting that the hope of the faithful is vain: and that is a foolish thing, to despise good things present and certain for an uncertain glory. There are innumerable others of the same sort, which come to the mind of man, and trouble and shake the faith of eternal life. The Lord, therefore, the faithful Physician of his people, lest they should feel any hindrance in this behalf, confirms these things strongly and in many ways: declaring the hope of the faithful to be most certain, and all things to be undoubted, which are or shall be taught of eternal life, of the felicity and glory of saints. And he places this assertion as it were in a bright ray, after he has certainly collected the sum of felicity, whereunto by and by he will add fuller things after the vision exhibited.

It is to be understood that the certainty of the blessed life is shown most clearly of all by these words: “And he said to me, write,” and so forth. Nevertheless, no weak reasons for the truth can be gathered from the preceding words either.

First, indeed, he says: “I, John, saw.” And we know John to be an apostle and witness of the truth, whose testimony it is unlawful to distrust. Saying therefore this godly man saw the things himself which he recounts; to doubt of them would be a wickedness.

Secondly, he hears a voice, and a great, resounding pronouncement coming from the throne, namely, from the twenty-four elders and angelic beings, and the entire heavenly host. And who can doubt their testimony, who already exist in everlasting blessedness? They know, and have experience of what felicity is; therefore they speak and testify that which is tried and known.

Moreover, he himself who sits on the throne speaks and testifies, saying: “Behold, I make all things new.” God is true, and in him there is no deceit. And saying this so plainly, he testifies that eternal life will be: and we see him declare also of what sort it will be. No place for doubtfulness is left hereafter.

And the things that he has shown and declared concerning the blessed life, he commands immediately to write. Things are written for a perpetual memorial of the thing, which we know to be true and substantial. For writings or testimonies which are written or made and sealed, by the law of nations and common custom of men, have the force of an undoubted testimony. But such letters or testimonies are made and sealed at the commandment of God. For God commands John to write those same things which are taught of the blessed life, and therefore they are true, undoubted, and infallible. As he himself immediately adds and says: “For these words are faithful and true, steadfast I say, and immutable.” What could be spoken more evidently than these? Here is also the authority of holy Scripture established. But he adds another thing almost more forceful: and he said to me, “It is done.” By this manner of speaking is signified either that the end has come, and all things are accomplished, like as it is used in chapter 16, or else that the thing which is spoken and believed to come, is so certain as though it were done already. We Germans, when we wish to signify that the thing which we have proposed, or promised and said, is sure, we are wont to say, \textit{Es ist gemacht}, “it is done.” Let us therefore believe assuredly these and all God’s words. Moreover, let us give our Lord God most hearty thanks, which with so great faith and diligence sustains and confirms our hope, and has commanded these mysteries of our salvation to be put in writing and published to the whole world in all ages. To him be glory forevermore. Amen.
